\documentclass{article}
\usepackage{amsmath,amssymb,amsthm}
\newtheorem{theorem}{Theorem}
\newtheorem{lemma}{Lemma}
\newtheorem{proposition}{Proposition}
\newtheorem{corollary}{Corollary}
\title{Real Analysis II}
\author{
\begin{tabular}{c}
Anurag (bmat2508@isibang.ac.in)\\[0.6em]
Instructor: C.R.E Raja
\end{tabular}
}
\date{}
\renewcommand{\thesubsection}{\arabic{subsection}}
\begin{document}
\maketitle
\section*{Class 1, January 9 2026}
\subsection{Riemann Integration}
The physical meaning of an integral of a function is the area under the curve with the x-axis. 
We intend to compute the said area by dividing the region into rectangular chunks.\\
Let $I=[a,b]$ be an interval. A partition $P$ of $I$ is a finite subset of $I$ that contains $a,b$.
For any partition $P$ there are $x_0,x_1,\dots,x_n$ such that \\$a=x_0<\dots<x_n=b$. We write $P=\{x_0,\dots,x_n\}$.
Let $f:I\to \mathbb{R}$ be a bounded function. Let $m,M$ be such that $m\leq f(x)\leq M$ for all $x\in I$. Let $P=\{x_0,\dots,x_n\}$ be a partition of $I$. 
Take $I_k=[x_{k-1},x_k]$ and call:
\begin{itemize}
\item $\Delta x_k=x_{k+1}-x_k$
\item $m_k=\inf_{x\in I_k}f(x)$
\item $M_k=\sup_{x\in I_k}f(x)$
\end{itemize}
Define $$\mathcal{L}(f,P)=\sum_{k=1}^{n}(\Delta x_k)m_k$$
$$\mathcal{U}(f,P)=\sum_{k=1}^{n}(\Delta x_k)M_k$$
Clearly, $\mathcal{L}(f,P)\leq \mathcal{U}(f,P)$. If $P$ and $Q$ are two partitions of $I$. We say that $Q$ is a refinement of $P$ if $P\subseteq Q$. We have the following result:
\begin{lemma}
In the above situation, $\mathcal{L}(f,P)\leq \mathcal{L}(f,Q)\leq \mathcal{U}(f,Q)\leq \mathcal{U}(f,P)$.
\end{lemma}
We will prove this later on.\\ \\ \\ \\ \\
Let $$\underline{\int_a^b}f(x)\text{ }dx=\sup_{P}\mathcal{L}(f,P)$$ be the Lower Riemann Sum and 
$$\overline{\int_a^b}f(x)\text{ }dx=\inf_{P}\mathcal{U}(f,P)$$
A bounded function $f$ in $I=[a,b]$ is called Riemann Integrable if the Lower Riemann Sum and the Upper Riemann sum are the same. In this 
case we write $f\in \mathcal{R}[a,b]$.\\
Let $P_1$ and $P_2$ be two partitions of $I$. Let $P=P_1 \cup P_2$. $P$ is a partition and its called a common refinement of $P_1$ and $P_2$.
\begin{proof}[Proof(Lemma 1)]
    Let $P=\{x_0<\dots<x_n\}$ and let $Q$ consist of just one more point $y$. Let $i$ be such that $x_{i-1}<y<x_i$.
    \\Let $$w_1=\inf f(x)\text{, }x_{i-1}\leq x\leq y$$
    and,
    $$w_2=\inf f(x)\text{, }y\leq x\leq x_i$$
    Clearly, we have $w_1\geq m_i$ and $w_2\geq m_i$, where $m_i=\inf f(x)\text{, }x\in [x_{i-1},x_i]$
Then $$\mathcal{L}(f,Q)-\mathcal{L}(f,P)=w_1(y-x_{i-1})+w_2(x_i-y)-m_i(x_i-x_{i-1})$$
$$=(w_1-m_i)(y-x_{i-1})+(w_2-m_i)(x_i-y)\geq 0$$
and hence the first part of the chain is proved. That $\mathcal{L}(f,Q)\leq\mathcal{U}(f,Q)$ is trivial. The last part can be proven analogously.
\end{proof}
\begin{proposition}
    $$\underline{\int_a^b}f(x)\text{ }dx\leq \overline{\int_a^b}f(x)\text{ }dx$$
\end{proposition}
\begin{proof}
Let $P_1$ and $P_2$ be two partitions with $P^{*}$ being their common refinement. Then we have:
$$\mathcal{L}(f,P_1)\leq\mathcal{L}(f,P^*)\leq\mathcal{U}(f,P^*)\leq\mathcal{U}(f,P_2)$$.
Keeping $P_2$ fixed and taking supremum over all $P_1$ yields
$$\underline{\int_a^b}f(x)\text{ }dx\leq \mathcal{U}(f,P_2)$$
finally taking infimum over all $P_2$ we get our desired result.
\end{proof}
\begin{proposition}
    $f\in \mathcal{R}[a,b]$ iff for each $\epsilon>0$ there is a partition $P$ such that 
    $$|\mathcal{U}(f,P^{*})-\mathcal{L}(f,P^{*})|<\epsilon$$
    for any refinement $P^{*}$ of P.
\end{proposition}
\begin{proof}
Assume first that $f$ is Riemann integrable over $I$. Let $\epsilon>0$. \\

$\int_a^b f=\inf$







For the converse, let us assume that for each $\epsilon>0$ there is a partition $P$ such that $$\mathcal{U}(f,P)-\mathcal{L}(f,P)<\epsilon$$
then 
$$\inf_{P}\mathcal{U}(f,P)-\sup_{P}\mathcal{L}(f,P)\leq \mathcal{U}(f,P)-\mathcal{L}(f,P)<\epsilon$$
and 
$$0\leq \overline{\int_a^b}f(x)\text{ }dx-\underline{\int_a^b}f(x)\text{ }dx<\epsilon$$
Let $\epsilon\to 0$. Then we get $$\overline{\int_a^b}f(x)\text{ }dx=\underline{\int_a^b}f(x)\text{ }dx$$ and thus $f\in \mathcal{R}[a,b]$.
\end{proof}
This is analogous to the Cauchy Criterion for sequences.




\section*{Class 2, January 13 2026}
\subsubsection{Properties}
\begin{proposition}
If $f,g\in \mathcal{R}[a,b]$ then $\alpha f+\beta g\in \mathcal{R}[a,b]$, for any $\alpha,\beta\in \mathbb{R}$.
\end{proposition}
\begin{proof}
$f\in \mathcal{R}[a,b]\implies \exists \text{ partitions } P_1,P_2 \text{ such that }$
$$\mathcal{U}(f,P_1)-\mathcal{L}(f,P_1)<\epsilon/2$$
and,
$$\mathcal{U}(g,P_2)-\mathcal{L}(g,P_2)<\epsilon/2$$
Let $P=P_1\cup P_2$. Then
$$\int_a^b f(x)\text{ }dx-\mathcal{L}(f,P)\leq \mathcal{U}(f,P)-\mathcal{L}(f,P)<\epsilon/2$$
and similarly,
$$\int_a^b g(x)\text{ }dx\leq \mathcal{U}(g,P)-\mathcal{L}(g,P)<\epsilon/2$$
We note that $$\mathcal{L}(f+g,P)\geq \mathcal{L}(f,P)+\mathcal{L}(g,P)$$ and $$\mathcal{U}(f+g,P)\leq \mathcal{U}(f,P)+\mathcal{U}(g,P)$$
And thus,
$$\int_a^b f(x)\text{ }dx+\int_a^b g(x)\text{ }dx-\mathcal{L}(f+g,P)\leq\int_a^b f(x)\text{ }dx+\int_a^b g(x)\text{ }dx-(\mathcal{L}(f,P)+\mathcal{L}(g,P))<\epsilon$$
We get a similar bound for the upper sum. This finally yields
$$\mathcal{U}(f+g,P)-\mathcal{L}(f+g,P)<2\epsilon$$.
Since $\epsilon$ was arbitrary, $f+g\in \mathcal{R}[a,b]$.
Now since 
$$\int_a^b f+g =\sup_{Q}\mathcal{L}(f+g,Q)$$
we have
$$\int_a^b f+\int_a^b g-\int_a^b (f+g)\leq \int_a^b f +\int_a^b g-\mathcal{L}(f+g,P)<\epsilon$$
letting $\epsilon\to 0$ we get 
$$\int_a^b f+\int_a^b g-\int_a^b (f+g)\leq 0$$
and similarly we have the reverse inequality. With this, we conclude that 
$$\int_a^b(f+g)=\int_a^b f+\int_a^b g$$
As for the scalar multiple case, let $\alpha\in \mathbb{R}$, $f\in\mathcal{R}[a,b]$.
Let $P=\{x_0,\dots,x_n\}$ be a partition.\\
\textbf{Case 1:} $\alpha\geq 0$. Then $$\mathcal{L}(\alpha f,P)=\sum \inf_{x\in I_k}\alpha f(x)\Delta x_k$$
$$=\sum \alpha \inf f(x)\Delta x_k=\alpha\mathcal{L}(f,P)$$
Similarly, $\mathcal{U}(\alpha f,P)=\alpha\mathcal{U}(f,P)$. Since this implies
$$\underline{\int_a^b} \alpha f=\alpha\underline{\int_a^b} f=\alpha\int_a^b f$$
and,
$$\overline{\int_a^b}\alpha f=\alpha\overline{\int_a^b} f=\alpha\int_a^b f$$
so $$\int_a^b \alpha f=\alpha\int_a^b f$$


\end{proof}
\begin{proposition}
    Let $f\in \mathcal{R}[a,b],\text{ }c\in [a,b]$. Then $f\in \mathcal{R}[a,c],f\in\mathcal{R}[c,b]$ and 
    $$\int_a^b f=\int_a^c f+\int_c^b f$$
\end{proposition}
\begin{proof}
Let $f_1=f|_{[a,c]}$ and $f_2=f|_{[c,b]}$. Now $f\in\mathcal{R}[a,b]\implies$ there exists a partition $P=\{x_0,\dots,x_n\}$ such that 
$$\mathcal{U}(f,P)-\mathcal{L}(f,P)<\epsilon$$
Let $c$ be such that $x_{i-1}<c<x_i$. Let $P^{*}=P\cup\{c\}$. Then 
$$\mathcal{U}(f,P^*)-\mathcal{L}(f,P^*)\leq \mathcal{U}(f,P)-\mathcal{L}(f,p)<\epsilon$$
Let $P_1=\{x_0,\dots,c\}$, $P_2=\{c,\dots,x_n\}$. With this, we have
$$\mathcal{U}(f,P^*)=\mathcal{U}(f_1,P_1)+\mathcal{U}(f_2,P_2)$$
and,
$$\mathcal{L}(f,P^*)=\mathcal{L}(f_1,P_1)+\mathcal{L}(f_2,P_2)$$
\end{proof}
\begin{proposition}
    Let $f\in\mathcal{R}[a,b]$. Then 
    \begin{enumerate}
    \item If $m\leq f(x)\leq M \text{ }\forall \text{ }x\in I=[a,b]$, then 
    $$m(a-b)\leq\int_a^b f\leq M(a-b)$$
    \item If $g\in\mathcal{R}[a,b]$ such that $g\leq f \text{ }\forall\text{ }x\in I$, then 
    $$\int_a^b g\leq\int_a^b f$$
    \item If $|f|\leq M$ then 
    $$|\int_a^b f|\leq M(b-a)$$
    \end{enumerate}
\end{proposition}
\begin{proof}
Suppose $m\leq f\leq M$ on $I$. Let $P=\{a,b\}$. Now 
$$\int_a^b f\geq \mathcal{L}(f,P)=(\inf_{I}f(x))(b-a)\geq m(b-a)$$
Similarly, 
$$\int_a^b f\leq \mathcal{U}(f,P)\leq M(b-a)$$
Suppose now that $g\in\mathcal{R}[a,b]$ and $g\leq f$. Then just consider the function $h=f-g\geq 0$. By the previous property, we have
$$\int_a^b(f-g)\geq 0$$
and by linearity of the integral, we have (2). The proof of the third property is left as an exercise.
\end{proof}
\begin{theorem}
Any continuous function on $[a,b]$ is Riemann Integrable.
\end{theorem}
\begin{proof}
Let $f:[a,b]\to\mathbb{R}$ be continuous. Then $f$ is bounded and $f$ is uniformly continuous. Let $\epsilon>0$ Then there exists some $\delta>0$ such that 
$$|f(x)-f(y)|<\epsilon$$
whenever $|x-y|<\delta$. Let $P=\{x_0<\dots<x_n\}$ be a partition of $[a,b]$ such that
$$\inf_{x\in I_k}f(x)=f(a_k)$$
and $$\sup_{x\in I_k}f(x)=f(b_k)$$
for some $a_k,b_k\in I_k$ as f is continuous. Then 
$$\mathcal{U}(f,P)-\mathcal{L}(f,P)=|\sum_{k=1}^n(x_k-x_{k-1})(f(b_k)-f(a_k))|\leq\sum_{k=1}^n (x_k-x_{k-1})|f(b_k)-f(a_k)|$$ 
$$a_k,b_k\in[x_{k-1},x_k]\implies |b_k-a_k|<\delta\implies |f(b_k)-f(a_k)|<\epsilon/(b-a)$$
and thus 
$$\mathcal{U}(f,P)=\mathcal{L}(f,P)\leq \frac{\epsilon}{b-a}\sum(x_i-x_{i-1})=\epsilon$$
and hence $f\in\mathcal{R}[a,b]$.
\end{proof}
\begin{theorem}
Any monotone function on $[a,b]$ is Riemann Integrable.
\end{theorem}
\begin{proof}
For simplicity, we assume $f$ to be monotonically increasing. \\
Then $f(a)\leq f(x)\leq f(b) \text{ }\forall \text{ }x\in [a,b]$, and hence $f$ is bounded. \\
Let $P=\{x-_0<\dots<x_n\}$ be a partition such that $$\sup_{x\in I_k}f(x)\leq f(x_k)$$
and $$\inf_{x\in I_k}f(x)\geq f(x_{k-1})$$
$$\mathcal{U}(f,P)-\mathcal{L}(f,P)=\sum(x_k-x_{k-1})(\sup_{x\in I_k}f(x)-\inf_{x\in I_k}f(x))$$
$$\leq \sum (x_k-x_{k-1})(f(x_k)-f(x_{k-1}))<\frac{\epsilon}{f(b)-f(a)}\sum f(x_k)-f(x_{k-1})<\epsilon$$
\end{proof}
\textbf{Homework:} If $\int_a^b f=0,\text{ }f\geq 0$ and $f$ is continuous on $[a,b]$. Then $f\equiv 0 \text{ }\forall\text{ }x\in [a,b]$.
\begin{proposition}
Let $f\in\mathcal{R}[a,b]$, $m\leq f(x)\leq M$ on $[a,b]$ and $\phi :[m,M]\to\mathbb{R}$ be continuous. Then $\phi\circ f\in \mathcal{R}[a,b]$.
\end{proposition}
\begin{proof}
Let $\epsilon>0$. Then there exists a $\delta >0$ such that 
$$|\phi(x)-\phi(y)|<\epsilon$$
whenever $|x-y|<\delta$. Now $f\in\mathcal{R}[a,b]\implies$ there exists a partition $P$ such that $\mathcal{U}(f,P)-\mathcal{L}(f,P)<\epsilon$. Let 
$Q$ be a refinement of $P$ so that the length of each of the subintervals in $Q$ is $\delta/4k$ where $k>1$ and $k>\phi(x)$ for all $x\in[m,M]$.
Now, $\mathcal{U}(f,Q)-\mathcal{L}(f,Q)<\epsilon$ and 
$$\mathcal{U}(\phi\circ f,Q)-\mathcal{L}(\phi\circ f,Q)=\sum[\sup \phi(f(x))-\inf \phi(f(x))](x_k-x_{k-1})$$
$$=\sum (M_k - m_k)(x_k-x_{k-1})\text{, where }M_k=\sup_{x\in I_k}\phi(f(x))\text{ and }m_k=\inf_{x\in I_k}\phi(f(x))$$
Let $a_k,b_k$ be the supremum and the infimum defined over the $k$-th subintervals respectively. 


\end{proof}
\section*{Class 6, January 27 2026}
\subsection*{Dirichilet's Test}
Let $f,g:[a,\infty)\to\mathbb{R}$ for some $a\in\mathbb{R}$. Assume that $g$ is monotonic and $\lim_{x\to\infty}g(x)=0$. If $f$ is continuous and if there is a positive $M$ such that 
$$\big{|}\int_a^x f(t)\text{ }dt\big{|}\leq M$$ for all $x>a$ and $g(x)$ is differentiable. Then $\int_a^{\infty} f(x)g(x)\text{ }dx$ converges.
\begin{proof}
    We will use the fact that $g$ is differentiable. We let 
    $$F(x)=\int_a^xf(t)dt$$ for all $x\geq a$. Then $|F(x)|\leq M$ for all $x\geq a$. Now
    $$\int_b^c f(x)g(x)dx=F(c)g(c)-F(b)g(b)-\int_b^c F(x)g'(x)dx\text{ (using integration by parts)}$$
As $F$ is bounded, 
$$F(c)g(c),F(b)g(b)\to 0$$ as $b\to\infty$. Then 
$$\big{|}\int_b^c F(x)g'(x)dx\big{|}\leq \int_b^c |F(x)|\cdot|g(x)|dx\leq M\int_b^c |g'(x)|dx $$
Since $g$ is monotonic, $g'$ does not change sign and thus we have 
$$\int_b^c |g'(x)|dx=|\int_b^c g'(x)dx|=|g(c)-g(b)|$$
where the last term goes to 0 as $b\to \infty$.
So the integral converges. 
\end{proof}
\begin{proposition}
    Let $f:[a,b]\to\mathbb{R}$ be a bounded function. Suppose $f\in\mathcal{R}[c,b]$ for all $c\in(a,b)$. Then $f\in\mathcal{R}[a,b]$ and $$\int_a^b f(x)\text{ }dx=\lim_{c\to a}\int_c^b f(x)\text{ }dx$$
\end{proposition}
\begin{proof}
    Let $\epsilon>0$. Let $M>|f(x)|$ for all $x\in[a,b]$. $f\in\mathcal{R}[a+\frac{\epsilon}{4M},b]$. There is a parition $P$ of this interval such that
    $$\mathcal{U}(P,f)-\mathcal{L}(P,f)<\frac{\epsilon}{2}$$
    Let $\tilde{P}=P\cup\{a\}$. Then 
    $$\mathcal{U}(\tilde{P},f)-\mathcal{L}(\tilde{P},f)$$
    $$=\big{[}\sup_{a\leq x\leq a+\frac{\epsilon}{4M}}f(x)-\inf_{a\leq x\leq a+\frac{\epsilon}{4M}}f(x)\big{]}+\mathcal{U}(P,f)-\mathcal{L}(P,f)$$
Now,
$$|\sup f(x)|\leq \sup |f(x)|\leq M\text{ for all }x$$
\end{proof}
\begin{theorem}
    Let $f:[a,b]\to\mathbb{R}$ be a bounded function. Suppose $f$ is continuous on $[a,b]$ except possibly at finitely many points. Then $f\in\mathcal{R}[a,b]$.
\end{theorem}
\begin{proof}
    Suppose $f$ is discontinuous on the set $F$ and for simplicity assume that $F=\{a\}$. The idea is to break the entire interval at points of discontinuity. Then it is enough to show that $f$ is Riemann integrable on an interval even if it is discontinuous at the endpoints.   
\end{proof}
We know that if $f$ is bounded on $[a,b]$ and $\int_c^b f(x)\text{ }dx$ exists for any $c\in[a,b]$ then this integral converges to $\int_a^b f(x)\text{ }dx$. But what if $f$ is not bounded on the interval?
\\
Let $f:(a,b]\to\mathbb{R}$ and $f\in\mathcal{R}[x,b]$ for every $x\in[a,b]$. We say that 
$$\int_a^b f(x)\text{ }dx$$
exists( or converges) if 
$$\lim_{x\to a}\int_x^b f(x)\text{ }dx$$ 
exists.\\
\textbf{Example:} $$\int_\epsilon^b x^{-p}\text{ }dx=\frac{b^{1-p}-\epsilon^{1-p}}{1-p}$$ for $p\neq 1$.
Then the integral converges to $\frac{b^{1-p}}{1-p}$ as $\epsilon\to 0$. \\
Similarly, we can define $\int_a^b f(x)$
\section*{Class 7, February 3 2026}
$f_n$ is said to converge pointwise to $f$ in $E$ if for every $\epsilon>0$ and and every $x\in E$ there exists $N\in \mathbb{N}$ such that for all $n\geq N$, 
$$|f_n(x)-f(x)|<\epsilon$$
$f_n$ is said to converge uniformly to $f$ if the same $N$ works for every $x\in E$.\\
We consider double sequences. Suppose $a_{m,n}=\frac{n}{n+m}$. Then $\lim_{n\to \infty}a_{m,n}=1$ and $\lim_{m\to\infty}a_{m,n}=0$. This shows that in general we cannot swaps double limits over a doubly indexed sequence. We cannot swap limits in pointwise convergence. \\
Take the example, $$f_m(x)=\lim_{n\to\infty}(\cos{m!\pi x})^{2n}=
\begin{cases}
0 \text{ if }x\in \mathbb{Q}\\
1 \text{ if }x\in \mathbb{Q}, \text{ }x=\frac{p}{q},\text{ }m\geq q
\end{cases}
$$
Taking $m\to\infty$, we find
$$f(x)=
\begin{cases}
0 \text{ if }x\notin \mathbb{Q}\\
1 \text{ if }x\in \mathbb{Q}
\end{cases}
$$
Next, we consider the example 
$$f_n (x)=n^2 x(1-x^2)^n\text{ , }x\in [0,1]$$
Then, 
$$\lim_{n\to\infty}f_n (x)=0$$
\subsection*{Uniform Convergence}
Let $E$ be a set and $f_n$ be a sequence of $\mathbb{R}/\mathbb{C}$ valued functions on $E$. We say that $f_n$ converges to uniformly to a function $f$ on $E$ if for each $\epsilon>0$, there is an $N\in\mathbb{N}$ such that for all $n\geq N$, $$|f_n(x)-f(x)|<\epsilon$$
for all $x\in E$. In this case we write $f_n\xrightarrow{u} f$ on $E$. Clearly, for a finite set $E$, pointwise and uniform convergence are all the same since we are able to get a maximum $N$ that works for all $x\in E$.
\begin{proposition}
If $f_n\to f$ on E, then for each $x\in E$, $f_n\to f$ pointwise in E.
\end{proposition}
We have a Cauchy Criterion here as well.
\begin{theorem}
    Let $\{f_n\}_{n\in\mathbb{N}}$ be a sequence of functions on $E$. Then $\{f_n\}_{n\in\mathbb{N}}$ converges uniformly to a function $f$ on $E$ iff for each $\epsilon>0$, there exists an $N\in\mathbb{N}$, $$|f_m(x)-f_n(x)|<\epsilon$$ for every $m,n\geq N$
\end{theorem}
\begin{proof}
    $(\impliedby)$
    Suppose $\{f_n\}_{n\in\mathbb{N}}$ converges uniformly to a function $f$ on $E$. Let $\epsilon>0$. Then by uniform convergence, 
    $$|f_n(x)-f(x)|<\epsilon$$
    for every $x\in E$ and for all $n\geq N$ for some $N\in\mathbb{N}$ depending on our choice of $\epsilon$. Then for any other $m\geq N$, we have by the triangle inequality,
    $$|f_m(x)-f_n(x)|<2\epsilon$$
    Since $\epsilon$ was arbitrary, we have our result.\\
    $(\implies)$Suppose to each $\epsilon>0$, there is an $N\in\mathbb{N}$ such that $$|f_n(x)-f_m(x)|<\epsilon$$ for all $m,n\geq N$ and for all $x\in E$. Thus the sequence $\{f_n(x)\}_{n\in\mathbb{N}}$ is Cauchy for every $x$. Then $\{f_n(x)\}_{n\in\mathbb{N}}$ converges for every $x\in E$. Then we let $m\to\infty$, to get
    $$-\epsilon\leq f_n(x)-f(x)\leq\epsilon$$
    But since $\epsilon$ was arbitrary, we have uniform convergence.
\end{proof}
\begin{theorem}[$L^{\infty}$ Test]
    Let $\{f_n\}_{n\in\mathbb{N}}$ be a sequence of functions on a set E and $f$ be a function on $E$. Then $\sup_{x\in E}{|f_n(x)-f(x)|}\to 0$ iff $f_n\xrightarrow{u}f$.   
\end{theorem}
\begin{proof}
    Let $f_n\xrightarrow{u}f$. So to each $e\epsilon$ there is an $N\in\mathbb{N}$ such that $|f_n(x)-f(x)|<\epsilon$ for all $n\geq N$ and $x\in E$. \\
     $(\iff)$ to each $\epsilon>0$, there is an $N\in\mathbb{N}$ such that
    $$\sup_{x\in E}|f_n(x)-f(x)|\leq \epsilon$$ for all $x\in E$.\\
    $(\iff)$ $$\sup_{x\in E}|f_n(x)-f(x)|\to 0 \text{ as }n\to \infty$$
\end{proof}
\textbf{Exercise: }Show that if $f_n\xrightarrow{u}f$, then $|f_n|\xrightarrow{u}|f|$.\\
\textbf{Exercise: }If $f_n\xrightarrow{u}f$ then does $\sup_{x\in E}f_n\to\sup_{x\in E}f$ ?
\subsection*{Series}
Let $\{f_n\}_{n\in\mathbb{N}}$ be a sequence of functions on a set $E$. We say that $\sum_{n=1}^{\infty}f_n$ converges uniformly on $E$ if the sequence of partial sums converges.
\begin{theorem}[Weirstrass M-Test]
    Let $\{f_n\}_{n\in\mathbb{N}}$ be a sequence of functions and let $M_n=\sup_{x\in E}|f_n(x)|$. Then $\sum_{n=1}^{\infty}f_n$ converges uniformly if $\sum_{n=1}^{\infty}M_n<\infty$. 
\end{theorem}
\begin{proof}
    Since $\sum_{n=1}^{\infty}M_n$ converges, for each $\epsilon$ there is an $N$ such that $\sum_{k=n}^{m}M_k<\epsilon$ for all $m\geq n\geq N$. Then 
    $$|\sum_{k=n}^{m}f_k|\leq |\sum_{k=n}^{m}M_k|<\epsilon$$
    Therefore, $\sum_{n=1}^{\infty}f_n(x)$ converges. 
\end{proof}
The M-Test works only in one direction.\\
\textbf{Example:} $f_n(x)=\frac{x^{2n}}{n^2}$ , $x\in [-1,1]$. Now $M_n=\sup |f_n(x)|=\frac{1}{n^2}$. Then since the sum of this series converges, the original series converges as well.\\
\textbf{Example:} $$f_n(x)=
\begin{cases}
0 \text{ if }x\neq \frac{1}{n}\\
\frac{1}{n} \text{ if }x=\frac{1}{n}
\end{cases}
$$
But we find that the sum of supremums does not converge. 
\begin{theorem}
    Let $f_n:[a,b]\to \mathbb{R}$ be a sequence of functions that converges uniformly to a function $f$ on $[a,b]$. 
If $\lim_{t\to x}f_n(x)=A_n$ exists, for some $x\in [a,b]$ for all $n\in\mathbb{N}$ then 
$$\lim_{n\to\infty}A_n$$
exists and equals 
$$\lim_{t\to x }f(t)=\lim_{t\to x}\lim_{n\to\infty}f_n(t)$$
Thus the limits can be exchanged.
\end{theorem}
\begin{proof}
    First we show that $\{A_n\}_{n\in\mathbb{N}}$ converges. For this we show that it is Cauchy. Let $\epsilon>0$, then there exists some $N$ such that $$|f_n(t)-f_m(t)|<\epsilon$$ for all $m,n\geq N$ and for all $t\in [a,b]$. Let $m,n\geq N$. Let $t\to x$ in the previous equation, then
    $$\lim_{t\to x}|f_m(t)-f_n(t)|\geq \epsilon$$ and since the absolute value is continuous, we have $$|A_m-A_n|\leq \epsilon$$. This implies that the sequence is Cauchy. Then $\{A_n\}_{n\in\mathbb{N}}$ converges.
Let $A=\lim A_n$. Then there is an $N$ such that $|f_n(t)-f(t)|<\frac{\epsilon}{3}$ for all $n\geq N$, $t\in[a,b]$ and $|A_n-A|<\frac{\epsilon}{3}$ for all $n\geq N$. Now $\lim f_n=A_n\implies$ there is a $\delta>0$ such that $|f_N(t)-A_N|$ for all $t\in (x-\delta,x+\delta)\cap [a,b]-\{x\}$. 
Finally,
$$|f(t)-A|\leq |f(t)-f_{N}(t)|+|f_{N}(t)-A_N|+|A_N-A|<\epsilon$$ for all $t\in(x-\delta,x+\delta)\cap [a,b]-\{x\}$.
From this, our result follows easily.
\end{proof}
\begin{corollary}
    Uniform limit of continuous functions is continuous.
\end{corollary}
\begin{proof}
    Let $\{f_n\}$ be a sequence of continuous functions on an interval $I$ such that $f_n\xrightarrow{u}f$. Let $x\in I$, for all $n$, we have
    $$\lim_{t\to x}f_n(t)-f_n(x)\implies \lim_{t\to x}f(t)=\lim_{n\to \infty}f_n(x)=f(x)\implies f\text{ is continuous at }x$$
\end{proof}
\section*{Class 8, February 6 2026}
\begin{corollary}
    Let $\{f_n\}$ be a sequence of functions in an interval $I$. Suppose $\sum_{n=1}^{\infty}f_n(x)$ converges uniformly and each of the $f_i(x)$ are continuous. Then their sum is also continuous.
\end{corollary}
\begin{proof}
    Let $s_n$ be the $n$-th partial sum. Then $s_n\to \sum_{n=1}^{\infty}f_n$. Now, the uniform limit of continuous function is continuous. Thus our result holds.
\end{proof}
\begin{theorem}
    There exists a continuous function that is nowhere differentiable.
\end{theorem}
\begin{proof}
    Let $f(x)=|x|$ defined on $[-1,1]$. Then $f$ is not differentiable just at one point. Now define $\phi:\mathbb{R}\to\mathbb{R}$ as 
    $$
    \phi(x)=
    \begin{cases}
    |x|, \text{ } |x|\leq 1\\
    |x+2k| ,\text{ }|x+2k|\leq 1\text{ where }k \text{ is an integer}
    \end{cases}
    $$
$\phi$ is the extension of $f$ with $\phi(x+2)=\phi(x)$. Therefore, $\phi$ is continuous and one can verify that $|\phi(x)-\phi(y)|\leq |x-y|$. Let $g(x)=\sum_{n=0}^{\infty}\big{(}\frac{3}{4}\big{)}^n \phi(4^n x)$
Now, 
$$|\big{(}\frac{3}{4}\big{)}^n \phi(4^n x)|\leq \big{(}\frac{3}{4}\big{)}^n$$
which leads to us conclude that $g$ converges uniformly and thus is continuous.
\\ Now let $x\in\mathbb{R}$. Fix $n\in\mathbb{N}$. Let $\delta_m =\pm \frac{1}{2}4^{m}$ where sign is chosen so that there are no integers between $4^m x$ and $4^m (x+\delta_m)$.
Let $\gamma_m=\frac{\phi(4^m (x+\delta_m))-\phi(4^m x)}{\delta_m}$
For $n>m$, $4^n \delta_m$ is an even integer and $\phi(4^n (x+\delta_m))=\phi(4^n x)$. Therefore, $\gamma_n=0$. 
Thus we have $$\sum_{n=0}^{\infty}\big{(}\frac{3}{4}\big{)}^n \frac{\phi(4^n (x+\delta_m))-\phi(4^n x)}{\delta_m}$$ 
One can conclude now that $g$ is not differentiable anywhere.     
\end{proof}
\begin{corollary}
    Let $\{f_n\}$ be a sequence of Riemann Integrable functions such that their sum converges uniformly on $[a,b]$. Then $\sum_{n}f_n \in \mathcal{R}[a,b]$ and $$\int_a^b \sum_{n}f_n=\sum_{n}\int_a^b f_n$$
\end{corollary}
\begin{theorem}
    Let $\{f_n\}$ be a sequence such that 
    \begin{enumerate}
        \item Each $f_n$ is continuous on $[a,b]$.
        \item $f_n$ converges to $f$ pointwise for all $x\in[a,b]$ for some continuous function $f$.
        \item $f_n$ is a decreasing sequence.
    \end{enumerate}
    Then $$\int_a^b f_n=\int_a^b f$$
    This is analogous to MCT. In fact, $f_n\xrightarrow{u}f$
\end{theorem}
\begin{proof}
    Assume first that $f$ is identically 0 on $[a,b]$. Let $g_n=f_n-f$. Then $g_n(x)\geq g_{n+1}(x)\text{ }\forall \text{ }x\in[a,b]$ and $g_n(x)\to 0$. Let $\epsilon>0$ and let$k_n=\{x \text{ : }g_n(x)>\epsilon \}$. Then we claim that some $k_n$ are empty. If this were false, assume that all $k_n$ were non-empty. Suppose $x_n\to x$. Then $x_n\in k_n \text{ }\forall\text{ }n$. Since $x_n\in k_n$, $g_n(x_n)>\epsilon$. Then letting $n\to\infty$ we find that $x\in k_n$ for all $n$. That means that $g_m(x)>\epsilon$ for all $m$ for some $x$. But this is a contradiction to the fact that $g_n$ converges to 0. Therefore, $k_n$ is empty for some $n$ and since this is a decreasing sequence, we conclude that $k_m$ is empty for all $m\geq n$ and all $x\in [a,b]$. Therefore, $g_n$ converges to 0 uniformly. 
 \end{proof}
\textbf{Example: }Let $f_n(x)=\frac{1}{nx+1}$. Note that $f_n$ is decreasing but does not converge unformly. This is fails since this is considered in an open interval. 

\end{document}
